\chapter{Teorema de la Aplicación Abierta}

\begin{coro}
    Una aplicación $f: X \longrightarrow Y$ es continua si y solo si la imagen inversa de cualquier abierto es un abierto. \\
    \noindent
    Es decir: si $V \subseteq Y$ es abierto $\implies f^{-1}(V) \subseteq X$ es abierto
\end{coro}

\begin{definicion}
    Sean $X, Y$ espacios normados y $T: X \to Y$ una aplicación lineal. Se dice que $T$ es \textbf{abierta} si $T(A)$ es abierto para todo $A \subseteq X$ abierto.
\end{definicion}

\begin{teo}[Teorema de la aplicación abierta]
    Sean $X, Y$ espacios de Banach y $T \in BL(X, Y)$ sobreyectiva, entonces $T$ es abierta.

\begin{proof}
    Para la demostración utilizamos la siguiente proposición: \\ \\
    \noindent
    \underline{Proposición.} Sean $X, Y$ espacios de Banach y $T \in BL(X, Y)$ sobreyectiva. Entonces 
    $$\forall r > 0, \exists s > 0 \text{ tal que } B^Y(0, s) \subseteq T(B^X(0, r))$$ \\
    \noindent
    Dada la proposición, sea ahora $G$ abierto en $X$ y sea $x_0 \in G$, tenemos 
    $$\exists r > 0 : B_r(x_0) \subseteq G \implies x_0 + B_r(0) \subseteq G$$
    Entonces $T(x_0) \in T(G)$ y, por la proposición 
    $$\exists s > 0 : B^Y(0, s) \subseteq T(B^X(0, r))$$
    Ahora
    \begin{align*}
        B^Y(T(x_0), s) & = T(x_0) + B^Y(0, s) \subseteq T(x_0) + T(B^X(0, r)) = T(x_0 + B^X(0, r)) \\
        &= T(B^X(x_0, r)) \subseteq T(G) \implies B^Y(T(x_0), s) \subseteq T(G) \quad \forall x_0 \in G
    \end{align*}
    entonces tenemos que $\forall y\in T(G), \exists s > 0 : B^Y(y, s) \subseteq T(G)$, es decir, $T(G)$ es abierto.
\end{proof}
\end{teo}

\begin{coro}
    Sea $T\in BL(X, Y)$ con $X, Y$ espacios de Banach y $T$ biyectiva, entonces
    $$T^{-1} \in BL(Y, X)$$
\begin{proof}
    Por el teorema de la aplicación abierta, $T$ es abierta, es decir, 
    \begin{equation*}
        \forall V\subseteq X, \text{ abierto} \implies T(V) \subseteq Y \text{ abierto}
    \end{equation*}
    y para demostrar que $T^{-1}$ es continua, basta con demostrar que la imagen inversa de un abierto por $T^{-1}$ es un abierto, es decir,
    \begin{equation*}
        \forall V \subseteq Y, \text{ abierto} \implies (T^{-1})^{-1}(V) \subseteq X \text{ abierto}
    \end{equation*}
    pero esto es cierto, ya que $(T^{-1})^{-1}(V) = T(V)$, y por ser $T$ abierta, $T(V)$ es abierto. \\ \\
    \noindent
    Por otra parte, la linealidad de $T^{-1}$ se sigue de la linealidad de $T$ y de la definición de inversa, es decir, para $y_1, y_2 \in Y$ y $\lambda \in \mathbb{R}$, tenemos
    \begin{align*}
        T^{-1}(\lambda y_1 + y_2) & = T^{-1}(\lambda T(T^{-1}(y_1)) + T(T^{-1}(y_2))) = T^{-1}(T(\lambda T^{-1}(y_1) + T^{-1}(y_2))) \\
        & = \lambda T^{-1}(y_1) + T^{-1}(y_2)
    \end{align*}
    y por lo tanto $T^{-1}$ es lineal, es decir, tenemos $T^{-1} \in BL(Y, X)$.
\end{proof}
\end{coro}

\begin{definicion}

    Sean $X, Y$ espacios topológicos, y $f$ una función de $X$ a $Y$; entonces, $f$ es un \tbf{homeomorfismo} si se cumple que:
    \begin{itemize}
        \item $f$ es una biyección
        \item $f$ es continua
        \item La inversa de $f$ es continua
    \end{itemize}
\end{definicion}

\subsection{Aplicaciones}
\noindent
Una de las aplicaciones más importantes del teorema de la aplicación abierta es la aproximación de soluciones de ecuaciones diferenciales. Queremos demostrar que si podemos aproximar una función ``difícil'' ($f$) mediante funciones ``fáciles'' (polinomios $p_n$), entonces las soluciones de la ecuación diferencial para esos polinomios ($x_n$) también aproximarán la solución real ($x$).  Vamos a desglosarlo paso a paso con rigor e intuición. \\ 

\noindent 
Consideremos el espacio 
$$C^1([0, 1]) = \{ x : [0, 1] \to \mathbb{R} \text{ continuas, derivables y } x' \text{ continua} \}$$ 
con la norma $\|x\|_{C^1} := \|x\|_{L^\infty} + \|x'\|_{L^\infty}$, el cual es de Banach. Sea 
$$X = \{ x \in C^1([0, 1]) \mid x(0) = x(1) = 0 \} \subseteq C^1([0, 1])$$ 
un subespacio cerrado en $C^1([0, 1])$ y, por lo tanto, de Banach. Sea 
\Func{T}{X}{C^0([0, 1])}{x}{\alpha x' + \beta x}
con $\alpha, \beta : [0, 1] \to \mathbb{R}$ continuas. Supongamos que $T$ sea biyectiva, es decir
\begin{equation*}
    \forall f \in C^0([0, 1]) \ \exists! x \in X : \alpha x' + \beta x = f \implies \begin{cases}
            \alpha x' + \beta x = f \\ x(0) = x(1) = 0 
    \end{cases}
    \text{ admite una única solución}
\end{equation*}

\begin{teo}[Teorema de Aproximación de Weierstrass]
    Si $f$ es una función real continua en $[a, b]$, entonces para cada $\epsilon > 0$, existe un polinomio $p$ tal que $|f(t) - p(t)| < \epsilon$ para todo $t \in [a, b]$. \\
    \noindent
    Esto equivale a decir, que el espacio de los polinomios es denso en $C^0([a, b])$ con la norma del supremo.
\end{teo}


\noindent 
Sean $(p_n) \in C^0([0, 1])$ polinomios tales que $p_n \to f$ uniformemente, es decir, existe la sucesión de polinomios $(p_n)$ tal que $\|p_n - f\|_{L^\infty} \to 0$, entonces tenemos que
\[ \forall p_n \ \exists! x_n : \begin{cases} \alpha x'_n + \beta x_n = p_n \\ x_n(0) = x_n(1) = 0 \end{cases} \text{ es decir, } T(x_n) = p_n \text{ con } (x_n) \in X. \text{ Se sigue que:} \]

\noindent 
Aplicando el corolario del teorema de la aplicación abierta, tenemos que $T$ es un homeomorfismo, es decir, $T$ es continua, biyectiva y su inversa es continua. Por lo tanto, como
\begin{equation*}
    p_n \to f \overset{T^{-1} \text{ cont.}}{\implies} T^{-1}(p_n) \to T^{-1}(f) \implies x_n \to x
\end{equation*}

\begin{ejercicio}
    Sean $X, Y$ espacios de Banach y $T \in BL(X, Y)$ inyectiva. Demostrar que 
    \begin{equation*}
        T^{-1}: T(X) \to X \text{ es continua} \iff T(X) \text{ es cerrado en } Y
    \end{equation*}
\begin{proof}
    Por hipótesis $T: X \to T(X)$ es biyectiva, al aplicar el teorema de la aplicación abierta y tener $T$ sobreyectiva (además de inyectiva por hipótesis del ejercicio), tenemos $T$ biyectiva, por lo tanto $T$ es un homeomorfismo.
    \begin{description}
        \item[$\Longrightarrow )$] Como $T(X) \simeq X$ y $X$ es de Banach, entonces $T(X)$ es de Banach, por lo tanto $T(X)$ es cerrado en $Y$, por ser subespacio de Banach, de $Y$.
        \item[$\Longleftarrow )$] Si $T(X)$ es cerrado en $Y$ de Banach, entonces $T(X)$ es de Banach, por lo tanto \mbox{$T \in BL(X, T(X))$} y por el teorema de la aplicación abierta, $T$ abierta y $T^{-1}:T(X)\to X$ es continua.
    \end{description}
\end{proof}
\end{ejercicio}

\begin{coro}\label{coro:equivalencia-normas}
    Sea $X$ un espacio vectorial y sean $(X, \|\cdot\|_1)$, $(X, \|\cdot\|_2)$ espacios de Banach, y supongamos que $\exists M > 0 : \|x\|_2 \leq M \cdot \|x\|_1$, entonces
    \begin{equation*}
        (X, \|\cdot\|_1), (X, \|\cdot\|_2) \text{ son equivalentes}
    \end{equation*}
\begin{proof}
    Sea $Id$ la aplicación identidad entre ambos espacios de Banach, entonces como
    \begin{equation*}
        Id: (X, \|\cdot\|_1) \longrightarrow (X, \|\cdot\|_2) \text{ es acotada} \implies Id \text{ es continua, lineal y biyectiva}
    \end{equation*}
    lo cual es claro, pues la identidad es lineal y biyectiva, y por la desigualdad dada en la hipótesis
    \begin{equation*}
        \exists M > 0 : \|x\|_2 \leq M \cdot \|x\|_1
    \end{equation*}
    $Id$ es acotada, es decir, continua, por ser equivalente ambas propiedades en operadores lineales. \\
    \noindent
    Por aplicación abierta $Id^{-1}$ es continua, entonces acotada, es decir
    \begin{equation*}
        \exists C > 0 : \|x\|_1 \leq C \cdot \|x\|_2
    \end{equation*}
    por lo que $(X, \|\cdot\|_1)$, $(X, \|\cdot\|_2)$ son equivalentes.
\end{proof}
\end{coro}

\section{Teorema del Gráfico Cerrado}

\begin{definicion}
    Sea $T: X \to Y$ una aplicación entre espacios normados. Se dice que 
    \begin{equation*}
        \Gamma_T = \{ (x, T(x)) \mid x \in X \} \subseteq X \times Y
    \end{equation*}
    es el \tbf{gráfico} de $T$. \\
\end{definicion}

\begin{definicion}
    Sea $T: X \to Y$ con $X, Y$ espacios normados. Se dice que $T: X \to Y$ es \textbf{cerrado} si $\Gamma_T$ es cerrado en $X \times Y$. \\
\end{definicion}

\begin{prop}
    Sea $T: X \to Y$ una aplicación entre espacios normados, entonces
    \begin{equation*}
        \Gamma_T \text{ es cerrado en } X \times Y \iff \forall (x_n) \subseteq X : x_n \to x \text{ y } T(x_n) \to y, \text{ entonces } y = T(x)
    \end{equation*}
\begin{proof}
    \begin{description}
        \item[$\Longrightarrow )$] Sea $(x_n) \subseteq X : x_n \to x$ y $T(x_n) \to y$, entonces $(x_n, T(x_n)) \to (x, y)$ en $\Gamma_T$ cerrado, es decir, $(x, y) \in \Gamma_T$ de lo cual $y = T(x)$. 
        \item[$\Longleftarrow )$] Sea $(x_n, y_n) \in \Gamma_T$ tal que $(x_n, y_n) \to (x, y)$, entonces
        $$(x_n, y_n) \in \Gamma_T \text{ con } y_n = T(x_n)$$ 
        y por lo tanto $(x_n, T(x_n)) \to (x, y)$. Por hipótesis $y = T(x)$.
    \end{description}
\end{proof}
\end{prop}


\observacion{
    Si $T$ es continua, entonces $\Gamma_T$ es cerrado. De hecho
    \begin{equation*}
        \text{si } x_n \to x, \text{ entonces } T(x_n) \to T(x)
    \end{equation*}
    y por lo tanto $(x_n, T(x_n)) \to (x, T(x)) \quad \forall (x_n) \to x$. \\
}

\noindent 
Si $T$ no es continua, entonces no es cierto en general, veámos dos ejemplo
\begin{itemize}
    \item Sea la función 
    \Func{T}{\mathbb{R}}{\mathbb{R}}{x}{\begin{cases} 1/x & x \neq 0 \\ 0 & x = 0 \end{cases}}
    no es continua en 0, pero dado su gráfico
    \begin{equation*}
        \Gamma_T = \{ (x, T(x)) : x \in \mathbb{R} \} = \{ (x, 1/x) : x \neq 0 \} \cup \{ (0, 0) \}
    \end{equation*}
    podemos ver que es cerrado en $\mathbb{R}$ (notar que no es completo).
    \item Sea la función
    \Func{T}{(C^1([0, 1]), \|\cdot\|_{L^\infty})}{(C^0([0, 1]), \|\cdot\|_{L^\infty})}{x}{x'}
    es lineal, para ver que no es continua, tomemos el ejemplo de esta función 
    \Func{f_n}{[0, 1]}{\mathbb{R}}{t}{\sin(nt)}
    donde tenemos ciertos puntos a destacar
    \begin{itemize}
        \item La norma en el dominio ($X$) es 
        $$\|f_n\|_\infty = \sup_{t \in [0,1]} |\sin(nt)| = 1\quad  \forall n$$
        \item La imagen del operador es 
        $$T(f_n) = f_n'(t) = n \cos(nt)$$
        \item La norma en el codominio ($Y$) es 
        $$\|T(f_n)\|_\infty = \sup_{t \in [0,1]} |n \cos(nt)| = n$$
    \end{itemize}
    y ahora vemos que mientras la ``entrada'' $f_n$ tiene siempre tamaño 1, la ``salida'' $T(f_n)$ crece hasta el infinito cuando $n \to \infty$, y entonces tenemos que 
    \begin{equation*}
        \nexists M > 0 : \|T(f_n)\|_\infty=n \leq M \cdot \|f_n\|_\infty=1 \quad \forall n
    \end{equation*}
    por lo que el operador no es acotado y, por tanto, no es continuo. Sin embargo, $\Gamma_T$ es cerrado porque si 
    \begin{equation*}
    x_n \xrightarrow{\text{unif.}} x \text{ y } T(x_n) = x'_n \xrightarrow{\text{unif.}} y \text{ en } Y
    \end{equation*}
    entonces $y = T(x)$ (TFC). Notar que dada la norma $\|\cdot\|_{L^\infty}$ es por eso que la convergencia por definición de la norma es convergencia uniforme.
\end{itemize}

\begin{teo}[Gráfico cerrado]
    Sean $X, Y$ espacios de Banach, tenemos que si
    \begin{equation*}
        T: X \to Y \text{ es lineal y cerrada } \implies T \text{ es continua}
    \end{equation*}
\begin{proof}
    Sean $\|\cdot\|_1, \|\cdot\|_2$ normas definidas sobre $X$ de la siguiente forma
    \begin{equation*}
        \|x\|_1 := \|x\|_X + \|T(x)\|_Y, \quad \|x\|_2 = \|x\|_X
    \end{equation*}
    es claro que $\|\cdot\|_2$ es de Banach, por ser igual a $\|\cdot\|_X$, pero para el caso de $\|\cdot\|_1$ tenemos que comprobarlo. \\ \\
    \noindent
    Para ello primero vemos que si $(x_n)$ es de Cauchy en $(X, \|\cdot\|_1)$, entonces $(x_n)$ es de Cauchy en $\|\cdot\|_X$ y $(\|T(x_n)\|)_Y$, ya que por ser de Cauchy en $(X, \|\cdot\|_1)$, entonces
    \begin{equation*}
        \forall \epsilon > 0, \exists N \text{ tal que si } n, m > N, \text{ entonces } \|x_n - x_m\|_1= \|x_n - x_m\|_X+\|T(x_n) - T(x_m)\|_Y < \epsilon
    \end{equation*}
    con lo que concluimos sabiendo que ambos términos (las normas) son no negativos que
    \begin{align*}
        \|x_n - x_m\|_X &\leq \|x_n - x_m\|_1 < \epsilon \\
        \|T(x_n) - T(x_m)\|_Y &\leq \|x_n - x_m\|_1 < \epsilon 
    \end{align*}
    es decir, $(x_n)$ es de Cauchy en ambas normas. Por lo que, sea $(x_n)$ de Cauchy en $(X, \|\cdot\|_1)$, entonces $(x_n)$ es de Cauchy en ambas normas, y como son respectivamente espacios de Banach, son completos, es decir
    \begin{equation}\label{eq:teo-grafico-cerrado}
        \exists x \in X : \|x_n - x\|_X \to 0 \quad \text{y } y \in Y : \|T(x_n) - y\|_Y \to 0
    \end{equation}
    por lo que para demostrar que $\|\cdot\|_1$ es de Banach, basta con ver
    \begin{equation*}
        \|x_n - x\|_1 = \|x_n - x\|_X + \|T(x_n) - T(x)\|_Y \to 0
    \end{equation*}
    donde teniendo que el primer término converge a $0$ por la ecuación \ref{eq:teo-grafico-cerrado}, el segundo lo hace si vemos que $y=T(x)$. Por hipótesis $\Gamma_T$ es cerrrado y se cumple que 
    \begin{equation*}
        (x_n, T(x_n)) \in \Gamma_T \ \forall n \quad \text{y} \quad (x_n, T(x_n)) \to (x, y) \in \Gamma_T \implies y = T(x)
    \end{equation*}
    por lo que $(X,\|\cdot\|_1)$ es de Banach. \\ \\
    Ahora, como $\|x\|_2 \leq \|x\|_1$ y por el Corolario \ref{coro:equivalencia-normas} del teorema de la aplicación abierta, $\|\cdot\|_2, \|\cdot\|_1$ son equivalentes, por lo tanto
    \begin{equation*}
        \exists M > 0 : \|x\|_1 \leq M \|x\|_2 \iff \|x\|_X + \|T(x)\|_Y \leq M \|x\|_X \iff \|T(x)\|_Y \leq (M - 1) \|x\|_X
    \end{equation*}
    de lo cual $T$ es acotado, es decir, continua.
\end{proof}
\end{teo}

\begin{coro}
    Sea $T: X \to Y$ lineal con $X, Y$ espacios de Banach, tenemos que si 
    \begin{equation*}
        T \text{ es cerrada y inyectiva} \implies T^{-1} \text{ continua}
    \end{equation*}
\begin{proof}
    Puesto que $T$ es cerrada, entonces $T$ es continua por Teorema del Gráfico Cerrado, por lo que teniendo ahora $T$ inyectividad, tenemos que
    \begin{itemize}
        \item Si $T$ es sobreyectiva, tenemos que $Y$ es de Banach por hipótesis.
        \item Si $T$ no es sobreyectiva, entonces $T(X)$ es un subespacio de $Y$ que es de Banach por ser homeormofo a $X$, y por lo tanto cerrado en $Y$.
    \end{itemize}
    Ahora teniendo $T$ biyectiva y $T\in BL(X, T(X))$ o $T\in BL(X,Y)$ por el corolario del teorema de la aplicación abierta, tenemos que $T^{-1}$ continua.
\end{proof}
\end{coro}

\section{Ortogonalidad en espacios de Banach}

\noindent
En esta sección tomaremos $X$ un espacio de Banach, $M\subseteq X$ un subespacio de $X$ y $N\subseteq X^*$ un subespacio de $X^*$.

\begin{definicion}
    Se define el \tbf{ortogonal} de $M$ como el subespacio $M^{\perp} \subseteq X^*$ tal que
    \begin{equation*}
        M^{\perp} := \{ f \in X^* \mid f(x) = 0 \quad \forall x \in M \}
    \end{equation*}
    Se define el \tbf{ortogonal} de $N$ como el subespacio $N^{\perp} \subseteq X$ tal que
    \begin{equation*}        
        N^{\perp} := \{ x \in X \mid f(x) = 0 \quad \forall f \in N \}
    \end{equation*}
    $M^{\perp}$ y $N^{\perp}$ son llamados \tbf{aniquiladores}.
\end{definicion}

\observacion{
    $M^{\perp}$ y $N^{\perp}$ son \underline{subespacios cerrados} incluso si $M$ y $N$ no lo son.
}

\begin{prop}
    Sea $M\subseteq X, N \subseteq X^*$, se cumple:
    \begin{enumerate}
        \item[i)] $(M^{\perp})^{\perp} = \overline{M}$
        \item[ii)] $\overline{N} \subseteq (N^{\perp})^{\perp}$. Si $X$ es reflexivo, se cumple $(N^{\perp})^{\perp} = \overline{N}$
    \end{enumerate}

\begin{proof}
    \begin{description}
        \item i) Sea $x\in M$, entonces 
        $$f(x) = 0 \ \forall f \in M^\perp \quad \text{ y } \quad (M^\perp)^\perp = \{ x \in X \mid f(x) = 0 \ \forall f \in M^\perp \}= \bigcap_{f \in M^{\perp}} \text{Ker}(f)$$
        de lo cual $x \in (M^\perp)^\perp$ y por lo tanto $M \subseteq (M^\perp)^\perp$. Como $(M^\perp)^\perp$ es cerrado, entonces $\overline{M} \subseteq (M^\perp)^\perp$. \\ \\
        \noindent
        Para el caso de la inclusión contraria, supongamos por reducción al absurdo que $(M^\perp)^\perp \not\subseteq \overline{M}$, es decir, 
        $$\exists x_0 \in (M^\perp)^\perp : x_0 \notin \overline{M}$$
        Consideremos $A = \{x_0\}, B = \overline{M}$. Tenemos que $A$ es compacto porque todo conjunto finito de puntos es compacto, y $B$ es cerrado y convexo, donde ser convexo es conclusión clara de ser un subespacio vectorial
        \begin{equation*}
            \lambda x + (1 - \lambda) y \in \overline{M} \quad \forall x, y \in \overline{M}, \lambda \in [0, 1]
        \end{equation*}
        Aplicando ahora el teorema de Hahn-Banach geométrico, $\exists \tilde{f} \in X^*, \alpha, \beta \in \mathbb{R}$ tales que
        \begin{equation*}
            \tilde{f}(x) \leq \beta < \alpha \leq \tilde{f}(x_0) \quad \forall x \in B=\overline{M}
        \end{equation*}
        Pero como $\overline{M}$ es un subespacio vectorial, entonces si tuvieramos $x'\in \overline{M}$ tal que $\tilde{f}(x')>0$, entonces $\tilde{f}(\lambda x') = \lambda \tilde{f}(x') > 0$ para $\lambda > 0$, es decir, $\tilde{f}$ no estaría acotada y sería una contradicción, porque antes teníamos que
        \begin{equation*}
            \tilde{f}(x) < \alpha \quad \forall x \in M
        \end{equation*}
        entonces $\tilde{f}=0 \ \forall x \in M$, es decir, $\tilde{f}\in M^{\perp}$. Ahora tenemos 
        \begin{equation*}
            0 < \alpha \leq \tilde{f}(x_0) \implies \tilde{f}(x_0) >0
        \end{equation*}
        pero dijimos que $\tilde{f}(x_0) = 0$ porque $x_0 \in (M^\perp)^\perp$, lo cual es una contradicción, por lo que $(M^\perp)^\perp \subseteq \overline{M}$ y por lo tanto $(M^\perp)^\perp = \overline{M}$.
         \\ \\
        \item ii) HACER % //TODO: hacer demostración para ii) (no está en los apuntes)
    \end{description}
\end{proof}
\end{prop}


\begin{prop}
    Sean $M_1, M_2$ subespacios de $X$, entonces
    \begin{enumerate}
        \item $M_1 \cap M_2 = (M_1^\perp + M_2^\perp)^\perp$
        \item $M_1^\perp \cap M_2^\perp = (M_1 + M_2)^\perp$
        \item $\overline{M_1^\perp + M_2^\perp} \subseteq (M_1 \cap M_2)^\perp$
        \item $(M_1^\perp \cap M_2^\perp)^\perp = \overline{M_1 + M_2}$
    \end{enumerate}
\end{prop}


\section{Operadores no acotados}

\begin{definicion}
    Sean $X,Y$ espacios de Banach y sea $D(T)$ un subespacio denso de $X$, un operador lineal $T$ definido como sigue
    \begin{equation*}
        T: D(T) \to Y
    \end{equation*}
    se dice \textbf{no acotado}.
\end{definicion}

\observacion{
    La definición extiende la definición de operador acotado porque $D(T)$ no es cerrado en general y, por lo tanto, no es necesariamente de Banach. 
}

\begin{observacion}
    En general, $T$ no envía acotados en acotados.
\end{observacion}

\begin{definicion}
    Dado $T: D(T) \to Y$ no acotado, si 
    
    \begin{equation*}
        \exists M > 0 \text{ tal que } \|T(x)\|_Y \leq M \cdot \|x\|_X \quad \forall x \in D(T)
    \end{equation*}
    entonces $T$ se dice \textbf{acotado} en $D(T)$.
\end{definicion}

\begin{definicion}
    Sea $T: D(T) \to Y$ no acotado, definimos los subespacios.
    \begin{itemize}
        \item $\text{Ker}(T):= \{ x \in D(T) \mid T(x) = 0 \}$
        \item $R(T) := \{ y \in Y \mid y = T(x) \text{ para algún } x \in D(T) \}$ (\tbf{rango})
    \end{itemize}
\end{definicion}

\begin{definicion}
    Sea $T: D(T) \to Y$ no acotado. $T$ se dice \tbf{densamente definido} en $X$ si $D(T)$ es denso en $X$. \\
\end{definicion}

\begin{prop}
    Si $T$ es cerrado, entonces $\text{Ker}(T)$ es cerrado.
\begin{proof}
    Sea $(x_n) \subseteq \text{Ker}(T)$ tal que $x_n \to x$, entonces $(x_n, T(x_n)) \in \Gamma_T$ cerrado, es decir, $(x_n, 0) \to (x, 0)$ de lo cual $T(x) = 0$, es decir, $x \in \text{Ker}(T)$ y por lo tanto $\text{Ker}(T)$ es cerrado.
\end{proof}
\end{prop}

\begin{ejemplo}
    Sea la función
    \Func{T}{(C^1([0, 1]), \|\cdot\|_\infty)}{(C^0([0, 1]), \|\cdot\|_\infty)}{u}{u'}
    está densamente definido en $(C^0([0, 1]), \|\cdot\|_\infty)$ y es cerrado, veámos estas afirmaciones por separado:
    \begin{itemize}
        \item Sabemos que $C^1([0, 1]) \subseteq C^0([0, 1])$ y que $C^1([0, 1])$ es denso en $C^0([0, 1])$, de hecho, dado $f \in C^0([0, 1])$, por el teorema de Weierstrass existe una sucesión de polinomios $(p_n)$ tal que $p_n \to f$ uniformemente, es decir, $\|p_n - f\|_\infty \to 0$, y como todo polinomio es de clase $C^1$, entonces $p_n \in C^1([0, 1])$ y por lo tanto $C^1([0, 1])$ es denso en $C^0([0, 1])$.
        \item Para ver que $T$ es cerrado, tenemos que demostrar que $\Gamma_T$ es cerrado en $C^1([0, 1]) \times C^0([0, 1])$. \\ \\
        \noindent
        Supongamos una sucesión $(u_n) \subseteq C^1([0, 1])$ tal que $u_n \to u$ en $(C^0, \|\cdot\|_\infty)$ y $T(u_n) = u_n' \to v$ en $(C^0, \|\cdot\|_\infty)$, entonces el resultado del cálculo clásico nos dice que \\

        \noindent
        ``Si una sucesión de funciones derivables converge uniformemente a $u$, y sus derivadas convergen uniformemente a $v$, entonces la función límite $u$ es derivable y su derivada es exactamente $v$ ($u' = v$)'' \\
        
        \noindent
        lo que implica que $u$ es derivable y $u' = v$, es decir, $u \in C^1([0, 1])$ y $T(u) = v$, por lo que $\Gamma_T$ es cerrado.
        
        

    \end{itemize}
\end{ejemplo}

\begin{ejercicio}
    Sean $X, Y$ espacios de Banach y $T: D(T) \to Y$ lineal, densamente definido y acotado. Entonces $T$ se puede extender unívocamente a un operador $\tilde{T}: X \to Y$ acotado. \\ \\
    \noindent
    Por hipótesis $\overline{D(T)} = X$, por lo que 
    \begin{equation*}
        \forall x \in X \quad \exists (x_n) \in D(T) \text{ tal que } x_n \to x
    \end{equation*}
    Definimos:
    \Func{\tilde{T}}{X}{Y}{x}{\lim\limits_{n \to +\infty} T(x_n)}
    de donde, si $x \in D(T)$ entonces $\tilde{T}(x) = T(x)$, es decir, $\tilde{T}$ es una \underline{extensión de $T$ }. Ahora veamos que es acotada, para ello buscaremos primero la igualdad $\|\tilde{T}\| = \|T\|$, veámoslo
    \begin{itemize}
        \item Tenemos que
        $$\|\tilde{T}(x)\| = \|\lim_n T(x_n)\| = \lim_n \|T(x_n)\| \overset{(*)}{\leq} \lim_n \|T\| \cdot \|x_n\| = \|T\| \cdot \|x\| \Rightarrow \|\tilde{T}\| \le \|T\| \Rightarrow \tilde{T} \in BL(X,Y)$$
        donde en $(*)$ se ha usado la acotación de $T$. 
        \item $x \in D(T): \|T(x)\| = \|\tilde{T}(x)\| \le \|\tilde{T}\| \cdot \|x\| \Rightarrow \|T\| \le \|\tilde{T}\|$.
    \end{itemize}
    De lo cual $\|T\| = \|\tilde{T}\|$. Es importante notar que la acotación se ha probado en la primera parte de la desigualdad, pero hemos hecho el segundo desarrollo para ver que se da la igualdad de las normas. \\
\end{ejercicio}


\subsection{Operador Adjunto}

\noindent
A continuación trabajaremos con el operador $T: D(T) \to Y$ densamente definido en $X$ y con $X, Y$ espacios de Banach, y usaremos el \underline{operador adjunto} $T^*$ definido en
\begin{equation*}
    D(T^*) = \{ g \in Y^* \mid \exists C > 0 \text{ tal que } |\langle g, T(u) \rangle| \le C \cdot \|u\| \quad \forall u \in D(T) \}
\end{equation*}
donde $\langle g, T(u) \rangle = g(T(u))$. \\

\begin{prop}
    Si el operador $T$ es acotado, entonces
    \begin{equation*}
        D(T^*) = Y^*
    \end{equation*}
\begin{proof}
    \begin{equation*}
        \| g \circ T(u) \| \le \|g\| \cdot \|T\| \cdot \|u\| \quad \forall g \in Y^* \Rightarrow D(T^*) = Y^*
    \end{equation*}
    donde hemos usado que $g$ es acotada por $g\in Y^*$. \\
\end{proof}
\end{prop}

\noindent
Sea la siguiente función
\Func{f}{D(T)}{\mathbb{R}}{u}{\langle g, T(u) \rangle}
donde $g \in D(T^*)$, entonces $f$ es acotada, ya que 
\begin{equation*}
    |f(u)| = |\langle g, T(u) \rangle| \le C \cdot \|u\| \quad \forall u \in D(T)
\end{equation*}
Por hipótesis $D(T)$ es denso en $X$, para poder definir $T^*$, y además como $f$ es acotada, entonces $f$ puede ser extendida a $F \in X^*$, es decir
\begin{equation*}
    \exists F \in X^* : F(u) = f(u) = \langle g, T(u) \rangle \quad \forall u \in D(T)
\end{equation*}
Dada esta conclusión podemos definir correctamente el operador adjunto $T^*$ de $T$ de la siguiente manera, veámoslo. \\

\begin{definicion}
    Sea $T: D(T) \longrightarrow Y$ densamente definido en $X$ con $X, Y$ espacios de Banach. Definimos el \textbf{operador adjunto} $T^*$ de $T$ como el operador
    \Func{T^*}{D(T^*)}{X^*}{g}{T^*(g)=F}
    para el cual tenemos que 
    \begin{equation*}
        \langle T^*g, u \rangle = (T^*g)(u) = F(u) = f(u) = \langle g, T(u) \rangle \quad \forall u \in D(T), \forall g \in D(T^*) 
    \end{equation*}
    donde $D(T^*)$ es el subespacio de $Y^*$ definido por
    \begin{equation*}
        D(T^*):= \{ g \in Y^* \mid \exists C > 0 \text{ tal que } |\langle g, T(u) \rangle| \le C \cdot \|u\| \quad \forall u \in D(T) \}
    \end{equation*}
    donde $\langle g, T(u) \rangle = g(T(u))$. \\
\end{definicion}

\begin{ejercicio}
    Sea $X=L^2(\Omega)$, $\Omega = (0,1) \subseteq \mathbb{R}$ y sea $H^1_0(\Omega) \subseteq L^2(\Omega)$ donde 
    \begin{equation*}
        H^1_0(\Omega) = \{ f \in L^2(\Omega) \mid \partial_x f \in L^2(\Omega) \text{ y } Tr f = 0 \}
    \end{equation*}
    Demostrar que la siguiente función
    \Func{T}{H^1_0(\Omega)}{L^2(\Omega)}{u}{u'}
    es un operador lineal, densamente definido y no acotado. \\

    \noindent
    Antes de iniciar la resolución vamos a tener claro que tipo de funciones contiene $H_0^1(\Omega)$, para ello, dado $f \in H_0^1(\Omega)$, entonces $f$ es una función de cuadrado integrable ($f\in L^2(\Omega)$), es decir
    \begin{equation*}
        \int_\Omega |f(x)|^2 dx=\int_0^1 |f(x)|^2 dx < +\infty
    \end{equation*}
    y debe cumplir dos condiciones
    \begin{itemize}
        \item $\partial_x f \in L^2(\Omega)$, es decir, la derivada de $f$ también es de cuadrado integrable, lo que implica que $f$ es una función de clase $C^1$.
        \item $Tr f = 0$, es decir, la traza de $f$ es nula, lo que implica que $f(0) = f(1) = 0$.
    \end{itemize}
    Sabiendo esto, demostremos lo que se pide para $T$.
    \begin{description}
        \item[No acotado.] Tomemos la sucesión de funciones 
        $$u_n(x) = \sin(2\pi n x)$$
        entonces $u_n \in H^1_0(\Omega)$ porque $u_n$ es de clase $C^1$, su derivada es de cuadrado integrable y su traza es nula. Ahora, veamos que el operador no es acotado, para ello, calculemos la norma de $T(u_n)$ y la norma de $u_n$, donde $T(u_m) = u_m' = 2\pi m \cos(2\pi m x)$.
        \begin{align*}
            \|u_n\| &= \sqrt{\int_0^1 \sin^2(2\pi n x) dx} =\sqrt{1/2} \quad \forall n \\
            \|T(u_n)\| &= 2\pi n \sqrt{\int_0^1 \cos^2(2\pi n x) dx} = 2\pi n \sqrt{1/2}
        \end{align*}
        con los desarrollos siguientes
        \begin{align*}
            \int_0^1 \sin^2(2\pi n x) dx &\overset{(*)}{=} \int_0^1 \frac{1 - \cos(4\pi n x)}{2} dx= \frac{1}{2} \int_0^1 1 dx - \frac{1}{2} \int_0^1 \cos(4\pi n x) dx= \\
            &= \frac{1}{2} - \frac{1}{2} \cdot \left[ \frac{\sin(4\pi n x)}{4\pi n} \right]_0^1 = \frac{1}{2} - \frac{1}{2} \cdot \frac{\sin(4\pi n) - \sin(0)}{4\pi n} = \frac{1}{2} - \frac{1}{2} \cdot \frac{0 - 0}{4\pi n} = \frac{1}{2} \\
            \int_0^1 \cos^2(2\pi n x) dx &\overset{(**)}{=} \int_0^1 \frac{1 + \cos(4\pi n x)}{2} dx= \frac{1}{2} \int_0^1 1 dx + \frac{1}{2} \int_0^1 \cos(4\pi n x) dx= \\
            &= \frac{1}{2} + \frac{1}{2} \cdot \left[ \frac{\sin(4\pi n x)}{4\pi n} \right]_0^1 = \frac{1}{2} + \frac{1}{2} \cdot \frac{\sin(4\pi n) - \sin(0)}{4\pi n} = \frac{1}{2} + \frac{1}{2} \cdot \frac{0 - 0}{4\pi n} = \frac{1}{2}
        \end{align*}
        donde en $(*)$ se ha usado la identidad trigonométrica $\sin^2(\alpha) = \dfrac{1 - \cos(2\alpha)}{2}$, y en $(**)$ se ha usado la identidad trigonométrica $\cos^2(\alpha) = \dfrac{1 + \cos(2\alpha)}{2}$. Ahora, tenemos que 
        \begin{equation*}
            \|T\|= \sup_{n} \frac{\|T(u_n)\|}{\|u_n\|} = \sup_{n} \frac{2\pi n \sqrt{1/2}}{\sqrt{1/2}} = 2\pi n \xrightarrow{n \to +\infty} +\infty
        \end{equation*}
        es decir, $T$ es no acotado. \\

        % //TODO: queda demostrar la linealidad y ser densamente definido, pero no está en los apuntes
    \end{description}
    Para el caso de varias dimensiones $n>1$, con $\Omega \subseteq \mathbb{R}^n$ abierto acotado, el operador 
    \Func{T}{H^1_0(\Omega)}{L^2(\Omega)}{u}{\nabla u = \left( \dfrac{\partial u}{\partial x_1}, \dfrac{\partial u}{\partial x_2}, \dots, \dfrac{\partial u}{\partial x_n} \right)}
    hereda la falta de acotación del operador derivativo unidimensional.
\end{ejercicio}
